\cotefile{15}
\scanpage{103}
\mspage{IV.8}

$\psi^*\psi_*(\handwritten{\mathscr{A}}')$ to $\handwritten{\mathscr{A}}'$ corresponding to the identity homomorphism on
$\psi_*(\handwritten{\mathscr{A}}')$. We denote this map by $\sigma_{\handwritten{\mathscr{A}}'}:\psi^*\psi_*(\handwritten{\mathscr{A}}')\handwritten{\longrightarrow}\handwritten{\mathscr{A}}'$. If $\handwritten{\mathscr{A}}$ is a
sheaf of rings on $X$, there is a canonical ring homomorphism from
$\handwritten{\mathscr{A}}$ to $\psi_*\psi^*(\handwritten{\mathscr{A}})$ corresponding to the identity homomorphism on $\psi^*(\handwritten{\mathscr{A}})$.
We denote this map by $\rho_{\handwritten{\mathscr{A}}}:{}\handwritten{\longrightarrow}\psi_*\psi^*(\handwritten{\mathscr{A}})$.

Now suppose $\Psi:(Y,\handwritten{\mathscr{A}}')\handwritten{\longrightarrow}(X,\handwritten{\mathscr{A}})$ is a map of ringed spaces ;
i.e. $\Psi=(\psi,\theta)$ where $\psi$ is a continuous map from $Y$ to $X$ and $\theta$
is a homomorphism of sheaves of rings $\theta:\handwritten{\mathscr{A}}\handwritten{\longrightarrow}\psi_*(\handwritten{\mathscr{A}}')$. Then for
any $\handwritten{\mathscr{A}}'$-module $\handwritten{\mathscr{M}}'$, $\psi_*(\handwritten{\mathscr{M}}')$ can be made into an $\handwritten{\mathscr{A}}$-module by
means of the operation of $\theta$. When $\psi_*(\handwritten{\mathscr{M}}')$ is considered a module
over $\handwritten{\mathscr{A}}$ in this way, it is denoted $\Psi_*(\handwritten{\mathscr{M}}')$. We have the following
canonical functorial homomorphisms [8, 0 ; \S\ 4.2] :

\[
\begin{aligned}
h_* &: \Psi_*(\underline{\mathrm{Hom}}_{\handwritten{\mathscr{A}}'}(\handwritten{\mathscr{M}}',\handwritten{\mathscr{N}}'))\handwritten{\longrightarrow}\underline{\mathrm{Hom}}_{\handwritten{\mathscr{A}}}(\Psi_*(\handwritten{\mathscr{M}}'),\Psi_*(\handwritten{\mathscr{N}}'))\\
t_* &: \Psi_*(\handwritten{\mathscr{M}}')\otimes_{\handwritten{\mathscr{A}}} \Psi_*(\handwritten{\mathscr{N}}')\handwritten{\longrightarrow}\Psi_*(\handwritten{\mathscr{M}}'\otimes_{\handwritten{\mathscr{A}}'}\handwritten{\mathscr{N}}').
\end{aligned}
\]

On the other hand we have the map $\psi^*(\theta):\psi^*(\handwritten{\mathscr{A}})\handwritten{\longrightarrow}\psi^*\psi_*(\handwritten{\mathscr{A}}')$, and
so we have $\sigma_{\handwritten{\mathscr{A}}'}:\psi^*(\theta):\psi^*(\handwritten{\mathscr{A}})\handwritten{\longrightarrow}\handwritten{\mathscr{A}}'$. $\handwritten{\mathscr{A}}'$ is an algebra over $\psi^*(\handwritten{\mathscr{A}})$
by the means of this map, and if $\handwritten{\mathscr{M}}$ is an $\handwritten{\mathscr{A}}$-module, we define
the $\handwritten{\mathscr{A}}'$-module $\Psi^*(\handwritten{\mathscr{M}})$ to be $\psi^*(\handwritten{\mathscr{M}})\otimes_{\psi^*(\handwritten{\mathscr{A}})}\handwritten{\mathscr{A}}'$.

There is a canonical functorial isomorphism [8, 0; 4.33.1]

\[
t^*:\Psi^*(\handwritten{\mathscr{M}})\otimes_{\handwritten{\mathscr{A}}'} \Psi^*(\handwritten{\mathscr{N}})\qquad\Psi^*(\handwritten{\mathscr{M}}\otimes_{\handwritten{\mathscr{A}}}\handwritten{\mathscr{N}}).
\]

It can be shown that there is a natural isomorphism from
$\mathrm{Hom}_{\handwritten{\mathscr{A}}}(\Psi^*(\handwritten{\mathscr{M}}),\handwritten{\mathscr{M}}')$ to $\mathrm{Hom}_{\handwritten{\mathscr{A}}'}(\handwritten{\mathscr{M}},\Psi_*(\handwritten{\mathscr{M}}'))$ for all $\handwritten{\mathscr{A}}$-modules $\handwritten{\mathscr{M}}$
and $\handwritten{\mathscr{A}}'$-modules $\handwritten{\mathscr{M}}'$. By means of this isomorphism we can define
the $\handwritten{\mathscr{A}}$-homomorphism $\rho_{\handwritten{\mathscr{M}}}:\handwritten{\mathscr{M}}\handwritten{\longrightarrow}\Psi_*\Psi^*(\handwritten{\mathscr{M}})$ and the $\handwritten{\mathscr{A}}'$-homomorphism

% archive sheet number: 102
